% =====================================================================
%  EPD-SDIA-4112 -- Statistique Bayesienne -- Master 1 SDIA
%  CHAPITRE 1 : Introduction et fondements bayesiens   (50 slides, 8 h)
%  Fichier de frames : aucune dependance a \title (build maitre + build seul)
% =====================================================================
\section{Chapitre 1 --- Fondements bayésiens}

% ---------------------------------------------------------------- 1
\chapterslide{Chapitre 1}{Introduction et fondements bayésiens}{8 h --- 5 sections --- 50 slides}{%
\textbf{Objectifs} --- distinguer les deux interprétations de la probabilité ;
énoncer et appliquer le théorème de Bayes ; définir prior, vraisemblance,
posterior et loi prédictive ; écrire un premier modèle en PyMC ;
lire un intervalle de crédibilité à 95~\%.}


% ---------------------------------------------------------------- 2
\begin{frame}{Pourquoi le bayésien ?}
\begin{columns}[t]
\begin{column}{0.52\textwidth}
\begin{ebox}[Trois raisons décisives]
\begin{itemize}
  \item \textbf{Données rares} : en santé, finance ou industrie, on ne dispose
        pas de 10\,000 observations. Le prior compense.
  \item \textbf{Connaissance experte} : l'expert sait déjà quelque chose ;
        le bayésien lui donne un langage pour l'écrire.
  \item \textbf{Incertitude honnête} : on répond par
        « $P(\theta > 0{,}5 \given y) = 0{,}97$ », pas par un rejet binaire.
\end{itemize}
\end{ebox}
\end{column}
\begin{column}{0.44\textwidth}
\begin{kbox}[Applications en M1 SDIA]
\begin{itemize}
  \item Dépistage et essais cliniques
  \item A/B testing en ligne
  \item Filtrage de spam, détection de fraude
  \item Systèmes de recommandation
  \item Modèles hiérarchiques et petits effectifs
  \item Apprentissage bayésien profond (BNN)
\end{itemize}
\end{kbox}
\end{column}
\end{columns}
\vspace{1mm}
\centering
\begin{tcolorbox}[colback=byCyan!8,colframe=byCyan,width=.94\textwidth,arc=2mm]
\centering\small
\textbf{Fil rouge du cours :} transformer une croyance a priori en
\textbf{distribution de probabilité} sur le paramètre, puis la mettre à jour par les données.
\end{tcolorbox}
\end{frame}

% ---------------------------------------------------------------- 3
\begin{frame}{Le problème qui met tout le monde d'accord}
\begin{columns}[t]
\begin{column}{0.48\textwidth}
\begin{ebox}[Énoncé]
Un test de dépistage a
\begin{itemize}
  \item sensibilité $P(T^+ \given M) = 0{,}99$
  \item spécificité $P(T^- \given \bar M) = 0{,}95$
  \item maladie rare : prévalence $P(M) = 0{,}01$
\end{itemize}
\vspace{1mm}
\textbf{Question :} le test est positif. Quelle est la probabilité d'être malade ?
\vspace{1mm}

\centering\large\textbf{Réponse intuitive la plus fréquente : 95 \%}
\end{ebox}
\end{column}
\begin{column}{0.48\textwidth}
\begin{center}
\begin{tikzpicture}[scale=.82,every node/.style={font=\scriptsize},
   level 1/.style={sibling distance=30mm,level distance=15mm},
   level 2/.style={sibling distance=15mm,level distance=17mm}]
\node[draw,rounded corners,fill=byLight,inner sep=2pt] {1000 personnes}
  child {node[draw,rounded corners,fill=byRed!12] {Malades (10)}
    child {node[fill=byRed!25,draw,rounded corners] {$T^+$ : 9,9}}
    child {node[fill=byLight,draw,rounded corners] {$T^-$ : 0,1}}}
  child {node[draw,rounded corners,fill=byBlue!12] {Sains (990)}
    child {node[fill=byOrange!25,draw,rounded corners] {$T^+$ : 49,5}}
    child {node[fill=byLight,draw,rounded corners] {$T^-$ : 940,5}}};
\end{tikzpicture}
\end{center}
\vspace{-1mm}
\begin{solbox}[Le calcul qui fâche]
\centering
$P(M \given T^+) = \dfrac{9{,}9}{9{,}9 + 49{,}5} = \mathbf{16{,}7\ \%}$
\end{solbox}
\end{column}
\end{columns}
\end{frame}

% ---------------------------------------------------------------- 4
\begin{frame}{Quiz d'ouverture (Kahoot / Mentimeter, 3 min)}
\begin{qbox}[Question 1 -- Diagnostic médical]
Même test (sensibilité 99 \%, spécificité 95 \%). La prévalence passe de 1 \% à 20 \%.
Quelle est la nouvelle probabilité d'être malade après un test positif ?
\vspace{2mm}
\begin{enumerate}
  \item[(A)] environ 17 \% \quad \textit{(inchangé : le test est le même)}
  \item[(B)] environ 83 \% \quad \checkmark
  \item[(C)] environ 99 \%
  \item[(D)] impossible à dire sans plus de données
\end{enumerate}
\vspace{2mm}
\textcolor{byBlue}{\textbf{Votez avant de calculer.}} Puis :
$P(M \given T^+) = \frac{0{,}99 \times 0{,}20}{0{,}99 \times 0{,}20 + 0{,}05 \times 0{,}80} = 0{,}832$.
\end{qbox}
\vspace{1mm}
\begin{kbox}[Le message de la séance]
La réponse ne dépend pas seulement du test : elle dépend de ce qu'on
\textbf{croyait avant}. C'est toute la statistique bayésienne en une phrase.
\end{kbox}
\end{frame}

% ---------------------------------------------------------------- 5
\begin{frame}{Trois siècles d'une idée contrariée}
\begin{center}
\begin{tikzpicture}[xscale=1.0,yscale=1.0]
\draw[very thick,byDeep!30,-{Latex[length=3mm]}] (0,0) -- (13.2,0);
\foreach \x/\y/\n/\t in {
   0.6/1/{1763}/{Théorème de Bayes (posthume)},
   2.6/-1/{1774--1812}/{Laplace : inférence bayésienne},
   4.8/1/{1920s}/{Fisher : le fréquentisme s'impose},
   6.9/-1/{1939}/{Jeffreys : priors non informatifs},
   9.0/1/{1990s}/{MCMC : renaissance (Gelfand--Smith)},
   11.0/-1/{2015+}/{PyMC, Stan, ArviZ},
   12.2/1/{2020s}/{IA bayésienne, LLM incertains}}
{
  \fill[byBlue] (\x,0) circle (2.2pt);
  \draw[byBlue!45] (\x,0) -- (\x,0.30*\y+0.0);
  \node[anchor=south,align=center,font=\scriptsize\bfseries,text=byDeep]
       at (\x,0.36*\y+0.06) {\n};
  \node[anchor=north,align=center,font=\tiny,text=byGrey,text width=2.3cm]
       at (\x,0.32*\y-0.02) {\t};
}
\end{tikzpicture}
\end{center}
\vspace{2mm}
\begin{ebox}[Pourquoi si longtemps ?]
\begin{itemize}
  \item Le calcul bayésien exige des intégrales $\int p(\theta) p(y \given \theta)\,\mathrm{d}\theta$,
        souvent sans forme close.
  \item Les ordinateurs et les algorithmes MCMC (chapitre 5) ont levé ce verrou
        \textbf{à la fin des années 1980}.
  \item Aujourd'hui une régression bayésienne tient en cinq lignes de \py{pm.sample()}.
\end{itemize}
\end{ebox}
\end{frame}
